%==============================================================================
\chapter{Fundamentação Teórica e Trabalhos Relacionados}\label{estadoDaArte}
%==============================================================================

Para fundamentar o projeto de um agente inteligente eficiente e sensível a custos, este capítulo estrutura a revisão bibliográfica em três pilares conceituais basilares, adotando a diretriz de concisão analítica:

\begin{enumerate}
    \item \textbf{Roteamento Dinâmico de Modelos em Agentes Autônomos}: mecanismos orientados a agentes para seleção e alternância dinâmica entre modelos de linguagem heterogêneos, visando a otimização de custos e precisão;
    \item \textbf{Geração Aumentada por Recuperação (RAG)}: injeção seletiva de conhecimento contextual para mitigação de alucinações e contenção de inchaço na janela de \textit{tokens};
    \item \textbf{Arquiteturas de Serviços e Redes Orientadas a Agentes (ADN)}: princípios de infraestrutura distribuída onde agentes cooperam sobre nós heterogêneos e hardwares aceleradores de inferência.
\end{enumerate}

\section{Agentes Autônomos e Roteamento Dinâmico de Modelos}

Tradicionalmente, a adoção de LLMs em sistemas corporativos e assistentes virtuais considerava a seleção de um modelo estático durante a fase de implantação. Em tais arquiteturas convencionais, todas as consultas recebidas pelo agente são invariavelmente encaminhadas para o mesmo provedor de modelo, independentemente da complexidade inerente da solicitação. Como observado por \citeonline{wang2024survey}, essa homogeneidade introduz ineficiências graves: perguntas triviais consomem o mesmo orçamento computacional e financeiro que instruções matemáticas ou encadeamentos lógicos complexos.

Para solucionar essa assimetria, surgiram propostas de roteamento de modelos (\textit{model routing}). Inicialmente, esses roteadores operavam como classificadores estáticos unidirecionais baseados em classificadores textuais simples. Contudo, em ambientes agenticos onde o sistema toma ações consecutivas, o roteamento puramente estático padece de um déficit informacional crítico quanto ao desempenho de execução real.

Essa limitação é enfrentada no trabalho de \citeonline{lance2026agentasrouter}, que introduz o paradigma \textit{Agent-as-a-Router} (ACRouter). No modelo proposto pelos autores, a tomada de decisão sobre qual LLM invocar não é tratada como uma simples classificação isolada, mas como um processo interativo formalizado em um ciclo de Contexto-Ação-Feedback (C-A-F):

\begin{itemize}
    \item \textbf{Contexto}: Avaliação da instrução corrente, restrições orçamentárias de tokens e histórico de interações;
    \item \textbf{Ação}: Seleção da LLM mais adequada (por exemplo, um modelo compacto de alta velocidade para tarefas de rotina versus um modelo de fronteira para raciocínio analítico);
    \item \textbf{Feedback}: O agente inspeciona o resultado gerado por meio de um módulo verificador (\textit{Verifier}), calculando métricas de confiança ou detectando falhas para acumular experiência ancorada em execução (\textit{execution-grounded experience}).
\end{itemize}

O trabalho de \citeonline{lance2026agentasrouter} demonstra empiricamente que o roteamento agentico expande a fronteira de Pareto entre acurácia e custo de tokens, reduzindo drasticamente o arrependimento acumulado (\textit{cumulative regret}) em comparação ao uso contínuo de um único modelo de topo de linha. Esse princípio constitui a pedra angular do motor de decisão adotado neste trabalho.

\section{Geração Aumentada por Recuperação (RAG) como Estratégia de Eficiência}

Em arquiteturas convencionais de agentes inteligentes, é frequente a prática de injetar vastos fragmentos documentais, históricos completos de conversação ou bases brutas de conhecimento diretamente no cabeçalho do \textit{prompt} do modelo. Essa abordagem, denominada de contexto exaustivo, acarreta severos efeitos colaterais: eleva quadraticamente o custo das requisições (uma vez que os provedores de nuvem tarifam os tokens de entrada), amplia o tempo até o primeiro token (\textit{Time to First Token} - TTFT) e induz o modelo ao fenômeno de dispersão atencional, culminando em alucinações.

O paradigma de Geração Aumentada por Recuperação (\textit{Retrieval-Augmented Generation} - RAG), introduzido formalmente no trabalho clássico de \citeonline{lewis2020retrieval}, estabelece um método rigoroso para desacoplar a memória paramétrica do modelo da base de conhecimento factual. Conforme esquematizado pelos autores, o RAG combina dois componentes síncronos:

\begin{equation}
    p(y|x) = \sum_{z \in \text{top-}k} p(z|x) \, p(y|x, z)
\end{equation}

\noindent onde $x$ representa a consulta de entrada, $z$ consiste nos documentos ou trechos relevantes recuperados a partir de um índice denso vetorial com base em similaridade semântica (cosseno ou produto escalar sobre vetores de \textit{embeddings}), e $y$ constitui a resposta gerada.

Sob a ótica da Engenharia de Software e da economia de recursos de agentes autônomos, o RAG atua como um filtro cirúrgico: apenas os $k$ fragmentos estritamente necessários para a resolução da demanda são extraídos e apresentados ao modelo de linguagem. Dessa forma, além de garantir ancoragem factual e diminuir drasticamente alucinações em domínios específicos, o RAG reduz a quantidade de \textit{tokens} de entrada transmitidos para as LLMs em até ordens de magnitude.

\section{Arquiteturas de Serviços e Redes Orientadas a Agentes (ADN)}

A execução distribuída e eficiente de agentes autônomos requer um arcabouço estrutural que vá além de simples invocações de scripts ad-hoc. Conforme explorado por \citeonline{deng2025asc}, a transição dos serviços computacionais orientados a funções para o paradigma de Computação de Serviços Agenticos (\textit{Agentic Services Computing} - ASC) redefine serviços como agentes autônomos orientados a objetivos, dotados de ciclo de vida próprio, percepção de contexto e capacidade de colaboração inter-agente.

No âmbito de infraestrutura e redes inteligentes, \citeonline{rossi2026adn} formalizam a visão pioneira de Redes Definidas por Agentes (\textit{Agentic-Defined Networking} - ADN). No modelo ADN, os agentes de inteligência artificial são posicionados como entidades de primeira classe embutidas ao longo de todo o contínuo de dispositivos, borda (\textit{edge}) e nuvem. A arquitetura ADN preconiza que os agentes devem:

\begin{itemize}
    \item Perceber continuamente o estado dos nós de execução e os recursos disponíveis no ambiente;
    \item Manter crenças e raciocinar autonomamente sobre restrições operacionais e métricas de desempenho estabelecidas por operadores de sistema;
    \item Coordenar dinamicamente o fluxo de trabalho e o despacho de processamento de acordo com capacidades heterogêneas de nós locais e remotos.
\end{itemize}

O alinhamento conceitual com o ADN fornece a base arquitetural para a inclusão de hardwares aceleradores de inferência dedicados, tais como as Unidades de Processamento de Linguagem (LPUs) desenvolvidas pela Groq \cite{groq2024lpu}. Ao contrário de GPUs convencionais, que enfrentam gargalos de largura de banda de memória durante a fase auto-regressiva de geração token a token, as LPUs oferecem execução determinística de alto rendimento, ultrapassando 500 a 800 tokens por segundo em modelos abertos (e.g., Llama 3 8B).

Essa assimetria de infraestrutura consolida a proposta deste trabalho: sob a ótica arquitetural do ADN, tarefas leves de orquestração, triagem de intenções e injeção de RAG podem ser atribuídas localmente ou a nós de inferência rápida LPU (Groq) a custos residuais, acionando serviços de nuvem de fronteira (mais onerosos) apenas sob demanda verificada. Complementarmente, aspectos de governança e monitoramento de execução de ferramentas por agentes conectam-se aos estudos de validação e controle desenvolvidos por \citeonline{souza2024firewalls}.

\section{Síntese Conceitual da Integração}

A Tabela \ref{tab:pilares} sumariza o papel exercido por cada trabalho basilar na sustentação do framework proposto nesta monografia:

\begin{table}[!htb]
\centering
\caption{Contribuição dos Pilares Teóricos para o Trabalho Proposto}
\label{tab:pilares}
\begin{tabular}{|p{4.5cm}|p{3.5cm}|p{6.5cm}|}
\hline
\textbf{Pilar / Artigo Base} & \textbf{Domínio} & \textbf{Papel no Framework Proposto} \\ \hline
\citeonline{lance2026agentasrouter} & Roteamento Agentico & Fornece o ciclo Contexto-Ação-Feedback (C-A-F) e o módulo verificador para decidir quando transitar entre modelos. \\ \hline
\citeonline{lewis2020retrieval} & Recuperação Vetorial (RAG) & Garante a recuperação estrita de contexto factual, reduzindo alucinações e contendo a expansão de tokens de entrada. \\ \hline
\citeonline{rossi2026adn} & Arquitetura de Redes Agenticas (ADN) & Estrutura a coordenação do fluxo de serviços e o despacho hierárquico sobre hardwares heterogêneos (Groq LPU vs. Nuvem). \\ \hline
\end{tabular}
\end{table}

A articulação sinérgica desses três pilares viabiliza a construção de um agente autônomo com consumo econômico balanceado de tokens, resposta em tempo real e estabilidade operacional.